\section{Lecture scientifique, ancrage et relecture automatique}
Semantic Reader explore l'enrichissement interactif des documents scientifiques, y compris des PDF existants \cite{ref15}. SciDaSynth extrait des tables structurées pilotées par les requêtes et soutient leur validation interactive \cite{ref16}. Ces travaux traitent l'extraction comme un matériau que l'utilisateur inspecte et corrige dans l'interface.

ALCE évalue les générations citées selon la fluidité, l'exactitude et la qualité des citations \cite{ref17}. Dans notre cas, retrouver une phrase par son identifiant et vérifier qu'elle soutient l'étape sont deux contrôles distincts. Montrer l'extrait dans sa mise en page d'origine aide le lecteur à juger, mais ne prouve pas la relation proposée.

Liang et collègues rapportent que 57,4\,\% des participants de leur étude prospective trouvent les retours GPT-4 utiles ou très utiles \cite{ref18}. MARG répartit lecture et discussion entre instances spécialisées et rapporte une réduction des commentaires génériques \cite{ref19}. Ces résultats justifient d'attacher chaque suggestion à un passage précis ; ils ne disent rien de la qualité de nos propres suggestions.

Zhuang et collègues synthétisent méthodes et difficultés de la relecture automatisée \cite{ref20}. La page technique de Stanford Agentic Reviewer décrit un workflow qui recherche des travaux sur arXiv \cite{ref21}. Nous distinguons cette documentation d'ingénierie d'une évaluation indépendante. Notre prototype ne calcule pas de score global de qualité scientifique.
